\underline{引理} 1.10. \underline{设二次模 $(V,Q)$ 为两个二次模 $(V_1,Q_1)$ 和 $(V_2,Q_2)$ 之和, 其中 $V_i$ 局部自由、秩为 $2m_i$, 且 $Q_i$ 非退化. 设 $V_i=W_i\oplus W'_i$ 为 $V_i$ 分解为两个全奇异子空间的分解 $(i=1,2)$. 则图}

\begin{center}
\begin{tikzcd}[column sep=huge,row sep=large]
C(Q) \arrow[rr,dash,"\sim"] \arrow[d,dash,"\wr"]
  & {} & C(Q_1)\otimes' C(Q_2) \arrow[d,equal] \\
\operatorname{End}(\bigwedge(W_1\oplus W_2))
  \arrow[r,phantom,"\simeq"]
  & \operatorname{End}(\bigwedge W_1\otimes\bigwedge W_2)
  \arrow[r,phantom,"\simeq"]
  & \operatorname{End}(\bigwedge W_1)\otimes'
    \operatorname{End}(\bigwedge W_2)
\end{tikzcd}
\end{center}

\underline{交换.}

%%%%%%%%%% source scan idx 76 | folio 69 %%%%%%%%%%

验证留给读者. 由此引理立即推出

\begin{equation}\tag{1.10.1}
e(W_1\oplus W_2)
=e(W_1)e(W_2)+(1-e(W_1))(1-e(W_2)).
\end{equation}

1.11. 下列构造借自 Micali 和 Villamayor. 若 $p:Z\longrightarrow S$ 为 $S$ 的二重平展覆盖, 则可将 $Z$ 视为 $S$ 上的 $\mathbb Z/2$-挠子(torsor). 以 $*$ 表示挠子的加法; 对于 $Z$ 的截面 $s$, 设 $e(s)$ 为 $p_*\mathcal O_Z$ 中在 $s$ 上等于一、在其余部分等于零的幂等元. 沿用 1.10 的记号, 存在唯一同构

\begin{equation}\tag{1.11.1}
\Spec(Z(Q_1))*\Spec(Z(Q_2))\longrightarrow\Spec(Z(Q)),
\end{equation}

其底层加法态射为

\begin{equation*}
+:\Spec(Z(Q_1))\times_S\Spec(Z(Q_2))
\longrightarrow\Spec(Z(Q)),
\end{equation*}

且在任意基变换后满足

\begin{equation*}
e(s+t)=e(s)e(t)+(1-e(s))(1-e(t)).
\end{equation*}

若 $s(W)$ 为 $\Spec(Z(Q))$ 中满足 $e(s(W))=e(W)$ 的截面, 则公式 (1.10.1) 变为

\begin{equation*}
s(W_1\oplus W_2)=s(W_1)+s(W_2).
\end{equation*}

\underline{命题} 1.12. \underline{设 $A$ 为域, $V$ 为 $A$ 上维数为 $2m>0$ 的向量空间, $Q$ 为 $V$ 上的非退化二次型, $W_1,W_2$ 为 $V$ 中两个维数为 $m$ 的全迷向子空间. 则 $e(W_1)=e(W_2)$ 当且仅当 $\dim(W_1/(W_1\cap W_2))$ 为偶数.}

\footnote{印本命题说 $V$ ``在 $V$ 上''的维数为 $2m$; 同一句中引入的基域为 $A$.}

%%%%%%%%%% source scan idx 77 | folio 70 %%%%%%%%%%

(参见 Bourbaki, \emph{Algèbre}, 第~9 章, \S6, 习题~28.) 设 $e_1,\ldots,e_m,f_1,\ldots,f_m$ 为 $V$ 的一组基, 使得

\begin{enumerate}[label=\alph*),leftmargin=2.5em]
\item $Q(\sum_i x_i e_i+\sum_i y_i f_i)=\sum_i x_i y_i$;
\item $W_1$ 由 $e_i$ 生成;
\item $W_2$ 由 $(e_i)_{1\leq i\leq s}$ 和 $(f_i)_{s<i\leq m}$ 生成.
\end{enumerate}

则 $V$ 为 $V_i=Ae_i+Af_i$ 之和, 且

\begin{equation*}
W_\alpha=\bigoplus_i(W_\alpha\cap V_i),
\qquad \alpha=1,2.
\end{equation*}

所以加法公式 1.11 将证明归结为 $\dim V=2$ 的情形. 若 $W_1=W_2$, 断言显然成立. 否则, 若 $W_1=Ae$ 且 $W_2=Af$, 并有 $\Phi(e,f)=1$, 则 $e(W_1)=fe$ 且 $e(W_2)=ef=1-fe$.

\medskip
2. \underline{二次超曲面}

2.1. 设 $k$ 为代数闭域, $V$ 为 $k$ 上维数为 $n+2$ 的向量空间. $k$ 上维数为 $n$ 的光滑二次超曲面, 是指同构于某个 $k$ 上普通二次型 $Q$ 的方程 $Q=0$ 所定义的 $\mathbb P(V^*)$ 的子概形的 $k$-概形 $X$.

\underline{注} 2.2. 当 $n=0,1,2$ 时, $X$ 为 $k$ 上维数为 $n$ 的光滑二次超曲面当且仅当

\begin{equation*}
\begin{array}{rcl}
n=0 &:& X\simeq\Spec(k)\amalg\Spec(k),\\[1ex]
n=1 &:& X\simeq\mathbb P^1_k,\\[1ex]
n=2 &:& X\simeq\mathbb P^1_k\times\mathbb P^1_k.
\end{array}
\end{equation*}
